\documentclass[12pt]{article}
\usepackage{amsmath,amssymb}

\begin{document}
\section*{{2024 USAMO Problems/Problem 1}}
Source: AoPS Wiki

\subsection*{Solution}
We can start by verifying that $n=3$ and $n=4$ work by listing out the factors of $3!$ and $4!$ . We can also see that $n=5$ does not work because the terms $15, 20$ , and $24$ are consecutive factors of $5!$ . Also, $n=6$ does not work because the terms $6, 8$ , and $9$ appear consecutively in the factors of $6!$ . Note that if we have a prime number $p>n$ and an integer $k>p$ such that both $k$ and $k+1$ are factors of $n!$ , then the condition cannot be satisfied. If $n\geq7$ is odd, then $(2)(\frac{n-1}{2})(n-1)=n^2-2n+1$ is a factor of $n!$ . Also, $(n-2)(n)=n^2-2n$ is a factor of $n!$ . Since $2n<n^2-2n$ for all $n\geq7$ , we can use Bertrand's Postulate to show that there is at least one prime number $p$ such that $n<p<n^2-2n$ . Since we have two consecutive factors of $n!$ and a prime number between the smaller of these factors and $n$ , the condition will not be satisfied for all odd $n\geq7$ . If $n\geq8$ is even, then $(2)(\frac{n-2}{2})(n-2)=n^2-4n+4$ is a factor of $n!$ . Also, $(n-3)(n-1)=n^2-4n+3$ is a factor of $n!$ . Since $2n<n^2-4n+3$ for all $n\geq8$ , we can use Bertrand's Postulate again to show that there is at least one prime number $p$ such that $n<p<n^2-4n+3$ . Since we have two consecutive factors of $n!$ and a prime number between the smaller of these factors and $n$ , the condition will not be satisfied for all even $n\geq8$ . Therefore, the only numbers that work are $n=3$ and $n=4$ . ~alexanderruan As listed above, $n=3$ and $n=4$ work but $n=5$ and $n=6$ don't. Assume $n>6$ , let $p$ be the smallest positive integer that doesn't divide $n!$ . It is easy to see $p$ is the smallest prime greater than $n$ . It is easy to see that both $p-1$ and $p+1$ are divisors of $n!$ Let $q$ be the smallest positive integer greater than $p$ that is 2 mod 6. We see that either $q=p+3$ or $q=p+1$ but $q$ , $q+1$ , $q+2$ must all be divisors of $n!$ giving a difference of 1 after a difference of 2. Therefore, all $n>6$ fail.
~mathophobia

\end{document}
